\CatchFileBetweenTags{\AlphaInvCodataExperimentalVal}{constants.tex}{AlphaInvCodataExperimentalVal}
\CatchFileBetweenTags{\AlphaInvCodataSigma}{constants.tex}{AlphaInvCodataSigma}

\CatchFileBetweenTags{\AlphaInvMorelExperimentalVal}{constants.tex}{AlphaInvMorelExperimentalVal}
\CatchFileBetweenTags{\AlphaInvMorelSigma}{constants.tex}{AlphaInvMorelSigma}

\CatchFileBetweenTags{\AlphaInvParkerExperimentalVal}{constants.tex}{AlphaInvParkerExperimentalVal}
\CatchFileBetweenTags{\AlphaInvParkerSigma}{constants.tex}{AlphaInvParkerSigma}

\CatchFileBetweenTags{\AlphaInvFanExperimentalVal}{constants.tex}{AlphaInvFanExperimentalVal}
\CatchFileBetweenTags{\AlphaInvFanSigma}{constants.tex}{AlphaInvFanSigma}

\CatchFileBetweenTags{\VonKlitzingVal}{constants.tex}{VonKlitzingVal}
\CatchFileBetweenTags{\VonKlitzingRkExperimentalVal}{constants.tex}{VonKlitzingRkExperimentalVal}
\CatchFileBetweenTags{\VonKlitzingRkSigma}{constants.tex}{VonKlitzingRkSigma}

\section{Empirical Correspondences and Observed Constants}
\label{sec:empirical_correspondences}

% State that Z_geo is dimensionless. To validate against experiment we translate
The structural index $Z_{\text{geo}}$ is strictly dimensionless. To validate it against experiment, we translate the geometric bounds into macroscopic observables using the standard dimensional scaffolding ($\hbar$, $c$, $\epsilon_0$) that observers employ to measure flux. These are correspondence identities; they do not derive dimensional constants from the lattice. The primary identification is with the inverse fine-structure constant, $\alpha^{-1}(0) \equiv Z_{\text{geo}}$. Three further corollaries follow: the elementary charge as the quantized flux unit, the von Klitzing resistance as the channel quantization, and the geometric load limit as the vacuum breakdown threshold.

\subsection{The Fine-Structure Constant and Quantization Gap}
\label{sec:fine_structure}

% Establish the primary empirical anchor: Z_geo corresponds to alpha^-1(0).
The structural index $Z_{\text{geo}}$ quantifies the substrate's intrinsic dimensionless coupling strength. In the analytic continuum dual, a static bound of this magnitude is identified with the inverse fine-structure constant, $\alpha^{-1}(0)$, with standard renormalization-group running governing dynamic variation around this fixed infrared baseline.

% Present experimental tension (kinematic vs perturbative vs synthesis).
Recent experimental values (\Cref{tab:alpha_comparison}) of $\alpha^{-1}(0)$ do not form a simple consensus and reflect ongoing metrological tension. Morel (2020) and Parker (2018) measure $\alpha^{-1}(0)$ kinematically using photon recoil (Rubidium and Cesium, respectively). Fan (2023) determines $\alpha^{-1}(0)$ via the electron anomalous magnetic moment ($g-2$). That these independent measurements are mutually exclusive is an open problem in standard metrology. CODATA (2022) synthesizes these conflicting inputs and its uncertainty reflects this tension.

\begin{table}[ht]
\centering
\begin{tabular}{lcc}
\toprule
\textbf{Source} & \textbf{Value ($\alpha^{-1}(0)$)} & \textbf{Dev.} \\
\midrule
Morel  (2020) & \AlphaInvMorelExperimentalVal & $\AlphaInvMorelSigma \sigma$ \\
Fan (2023) & \AlphaInvFanExperimentalVal & $\AlphaInvFanSigma \sigma$\\
Parker (2018) & \AlphaInvParkerExperimentalVal & $\AlphaInvParkerSigma \sigma$ \\
CODATA (2022) & \AlphaInvCodataExperimentalVal & $\AlphaInvCodataSigma \sigma$ \\
\bottomrule
\end{tabular}
\caption{Comparison of the $E_8$-Persistence structural index $Z_{geo}$ against experimental determinations of $\alpha^{-1}(0)$. Deviation is computed as $|Z_{geo} - \alpha^{-1}(0)_{exp}| / \sigma_{exp}$.}
\label{tab:alpha_comparison}
\end{table}

% Explain Quantization Gap: continuum approximation saturates discrete hardware.
\paragraph{The QED extraction tension and the Quantization Gap.}
Standard QED serves as the analytic continuum dual to this discrete architecture. Dimensional regularization approximates the substrate at low loop orders by treating spacetime as continuous. However, because the substrate possesses finite channel capacity ($\nu=16$ per sector, $N=32$ total), this continuum approximation must eventually break down when its combinatorial complexity outstrips the discrete hardware. The number of QED diagrams grows factorially with loop order, while the available combinatorial state space of the substrate scales only exponentially as $\nu^n = 16^n$. Because factorial growth strictly overtakes finite exponential capacity, the continuous expansion must encounter the substrate's information-theoretic resolution limit at high loop orders. We term this regime the \textbf{Quantization Gap}.

% Connect Dyson's asymptotic observation to the structural mechanism.
Dyson's 1952 observation\cite{dyson_divergence_1952} that the QED perturbative series is asymptotic identified the mathematical signature; the $E_8$ substrate provides the structural mechanism. The continuum dual does not fail from a pathology in the coefficients, but from exhausting the finite Shannon capacity of the underlying geometric hardware.

% Distinguish kinematic vs perturbative measurements.
Morel (2020) measures $\alpha$ kinematically, a direct macroscopic probe of the geometric impedance that does not rely on quantum loop corrections. Fan (2023), by contrast, extracts $\alpha$ from a 10th-order (5-loop) QED perturbative expansion, placing it precisely in the regime where the continuum dual approaches capacity saturation. The $\AlphaInvFanSigma\sigma$ deviation observed in Fan (2023) is consistent with the expected onset of the Quantization Gap, though it does not alone constitute confirmation.

% State falsification criterion.
\paragraph{Falsification criterion.} If pure kinematic measurements continue to converge to $\alpha^{-1}(0) \approx \num{\AlphaInvVal}$ while 6+ loop QED extractions systematically diverge from this value, the Quantization Gap is confirmed. Conversely, convergence of both methods to a single identical value differing from our geometric prediction by $>5\sigma$ would decisively falsify the $E_8$-Persistence theory.

\subsection{The Geometric Origin of Elementary Charge}
\label{sec:elementary_charge}
% Frame Z_geo as flow constraint through chi=2 boundary. Identify emergent flux quantum with e.
On the $E_8$ substrate, $Z_{\text{geo}}$ acts as a strict flow constraint: it restricts the amount of information flux that can propagate through a $\chi=2$ topological boundary without violating structural coherence. In standard macroscopic physics, this emergent quantum of flux is observed as the elementary charge $e$. Applying the empirical dimensional isomorphism that maps $Z_{\text{geo}}$ to the inverse fine-structure coupling yields:

\begin{equation}
    e = \sqrt{\frac{4\pi\epsilon_0\hbar c}{Z_{geo}}}
\end{equation}

The square-root form reflects that entropic action scales quadratically with information flux, while the geometric $4\pi$ factor arises from the isotropic boundary measure of the spherical topological defect ($S^2$, $\chi=2$) enclosing the flux.

\subsection{Macroscopic Channel Resistance}
\label{sec:channel_resistance}
A geometric impedance should manifest as a measurable macroscopic transmission resistance. The theoretical resistance of a single discrete channel follows from three structural facts:

\begin{enumerate}
    \item \textbf{Single channel limit:} The vacuum substrate supports discrete transmission channels.
    
    \item \textbf{Spinor double cover:} Because the carrier states are spinors (\Cref{sec:chiralcapacity}), completing a closed circuit to return to the initial phase requires traversing the manifold twice ($720^\circ$). The measured resistance of a single loop is therefore the vacuum impedance shared across two windings: $R_{\text{channel}} = Z_{\text{channel}}/2$.
    
    \item \textbf{Geometric coupling:} The dimensionless scaling between a single internal channel and the bulk medium is $Z_{\text{channel}} / Z_0 = Z_{\text{geo}}$, where $Z_0$ is the characteristic impedance of free space.
\end{enumerate}

Combining these yields the structural prediction:
\begin{equation}
    R_{\text{channel}} = \frac{Z_{\text{channel}}}{2} = \frac{Z_0}{2} \, Z_{\text{geo}} \approx \mathbf{\VonKlitzingVal}.
\end{equation}
The factor of $1/2$ originates strictly from the spinor double cover; $Z_{\text{geo}}$ supplies the dimensionless geometric coupling. The observer imposes $Z_0 = \sqrt{\mu_0/\epsilon_0}$ as dimensional scaffolding to map this dimensionless geometric resistance into SI units (ohms).

\paragraph{Verification via the von Klitzing constant ($R_K$):}
Standard electrodynamics defines $R_K \equiv h/e^2 = Z_0 / 2\alpha$. Because $Z_{\text{geo}}$ corresponds to $\alpha^{-1}(0)$, the structural prediction $R_{\text{channel}}$ matches $R_K$ exactly. The quantum Hall effect measures this single-channel resistance with parts-per-billion precision, independently of high-energy $\alpha^{-1}$ extractions:

\begin{itemize}
    \item \textbf{Experimental Value (CODATA 2022):} \VonKlitzingRkExperimentalVal
    \item \textbf{Precision:} The $E_8$-Persistence geometric prediction matches the observed quantum Hall resistance to within \VonKlitzingRkSigma $\sigma$.
\end{itemize}
Matching $R_{\text{channel}}$ to $R_K$ demonstrates that the quantum Hall plateau resistance is the direct macroscopic signature of single-channel $E_8$ spinorial transport.


\textit{Physical consequence: the flatness of QHE plateaus.} 
The quantization $R = R_{\text{channel}}/n$, regardless of impurities, reflects the lattice's discrete bit-depth: resistance is quantized by channel count, offering a structural origin for the topological protection traditionally attributed to electron wavefunction topology (Chern numbers) \cite{thouless_quantized_1982, kohmoto_topological_1985}.

\CatchFileBetweenTags{\PlanckChargeAttenuationVal}{constants.tex}{PlanckChargeAttenuationVal}
\CatchFileBetweenTags{\PlanckChargeRatioPercentageVal}{constants.tex}{PlanckChargeRatioPercentageVal}

\subsection{The Geometric Load Limit}
\label{sec:load_limit}

The geometric impedance $Z_{\text{geo}}$ enforces a safety margin between operating signal and substrate breakdown. If the substrate can theoretically support a maximum unitary flux density $q_{\text{max}}$, the lattice impedance attenuates the stable, propagating flux $q_{\text{operating}}$ by the square root of the impedance:

\begin{equation}
    q_{operating} = \frac{q_{max}}{\sqrt{Z_{geo}}} \approx \frac{q_{max}}{\num{\PlanckChargeAttenuationVal}}
\end{equation}

\textbf{Verification via the Planck Charge Ratio:}
% Verify via Planck charge ratio. Frame as structural necessity.
In standard physics, the absolute capacity of the vacuum is represented by the Planck charge $q_P$, and the fundamental operating flux is the elementary charge $e$. The empirical ratio $e/q_P \approx 0.085$ has historically lacked a structural explanation. On the $E_8$ substrate, this attenuation is a structural necessity preventing immediate vacuum breakdown.

\paragraph{Physical consequence: Safe load limit.}
The electron represents the safe operating limit of the vacuum. While the substrate can theoretically support a unitary charge ($q_P$), the geometric impedance restricts propagating charge to $\approx \num{\PlanckChargeRatioPercentageVal}\%$ of this maximum to maintain structural coherence.

\paragraph{Physical consequence: Vacuum breakdown as capacity limit (Schwinger limit).}
In the macroscopic QED dual, the standard model utilizes the electron mass $m_e$ as a characteristic scale for the Schwinger critical field $E_c = m_e^2 c^3/e\hbar$, marking the threshold where external electric fields spontaneously produce $e^+e^-$ pairs \cite{schwinger_gauge_1951}. Like $\hbar$, $c$, and $\epsilon_0$, the mass $m_e$ is part of the dimensional scaffolding supplied by the observer's measurement framework; it is not derived from the lattice. The Coulomb field of an elementary charge, evaluated at the reduced Compton wavelength $\lambdabar = \hbar/m_e c$, scales with this critical value by exactly the fine-structure constant $\alpha$:
\begin{equation}
    \frac{e}{4\pi\epsilon_0 \lambdabar^2} = \alpha E_c = \frac{E_c}{Z_{\text{geo}}}
\end{equation}
The dimensionless geometric impedance factor $Z_{\text{geo}} \approx \num{\AlphaInvVal}$ therefore represents the safety factor between the characteristic field of an isolated fundamental charge and the vacuum breakdown threshold. Any localized excitation attempting to concentrate flux beyond this margin hits the non-linear Schwinger ceiling, resolving into particle-antiparticle pairs and enforcing the substrate's ultimate capacity limit.

All four empirical correspondences (the fine-structure constant, the elementary charge, the quantized channel resistance, and the vacuum breakdown threshold) are constrained by the same geometric impedance $Z_{\text{geo}}$ of the $E_8$ substrate. The dimensionless relations contain no continuous parameters; the dimensional scaffolding ($\hbar$, $c$, $\epsilon_0$, $m_e$) is supplied by the observer's measurement framework, not derived from the lattice.